\subsection*{Questão 05}

\textit{Qual a relação de período e frequência? }

\textbf{Resposta:} 
A relação entre eles é proporcionalidade inversa. Em termos simples: quanto mais rápido algo se repete, menor é o tempo que leva para completar uma única repetição.

\begin{itemize}
    \item \textbf{Período (T):} É o intervalo de tempo necessário para que um ciclo completo ocorra. Sua unidade no Sistema Internacional (SI) é o segundo (s).
    \item \textbf{Frequência ($f$):} É o número de ciclos que ocorrem em uma determinada unidade de tempo (geralmente 1 segundo).
    Sua unidade no SI é o Hertz (\hz), onde $1 \hz = 1$ ciclo/segundo.
\end{itemize}

\textbf{Fórmula Matemática:}
\begin{equation}
    f = \frac{1}{T} \qquad ou \qquad T = \frac{1}{f}
\end{equation}

Portanto, se o período aumenta, a frequência diminui, e vice-versa.
Por exemplo, se um ciclo leva 0,5 segundos para se completar (T = 0,5 s), a frequência será de 2 Hz ($f = \frac{1}{0,5} = 2$ Hz).
Se o período for de 0,25 segundos (T = 0,25 s), a frequência será de 4 Hz ($f = \frac{1}{0,25} = 4$ Hz).